\cotefile{53}
\scanpage{23}
% loc 1–2, opening bracket
[\unsure{cel|cela|cet|ceci} exprime que l'hom $A\mathrel{\unsure{\text{---}|\to}}k'$ est un hom d'algèbres]

% loc 3–6, Si G
Si $G$ est muni d'une loi de comp. de $=$ $\mathrm{Dist}(D)$, celle-ci est ass. (comm) \unsure{sss|ssi} celle de $G$ l'est, et \deleted{\unsure{unitaire}} $G$ admet une unité bilatère ssi $\mathrm{Dist}(D)$ a une unité \deleted{\unsure{élément}} $e$ telle que $\delta(e)=e\otimes e$.

% loc 7–9, Dire que
Dire que $G$ opère sur un module $M$ \inserted{\unsure{(à g.)}} signifie que $\mathrm{Dist}(G)$ y opère à gauche -- et cette opération étend celle de $G$.

% loc 10–14, NB
\uline{NB} Dans certains cas, $\mathrm{Dist}(\unsure{X|x})$ se reconstitue à partir de $\mathrm{dist}(X,k)=\unsure{\mathrm{Dst}|\mathrm{Dist}}(X)(k)$, p.ex. si $A$ est \deleted{\unsure{de t.f.}} projectif de t.f. sur $k$, ou plus gén. si $A$ est somme de modules projectifs de t.f (p.ex. \deleted{\unsure{si}} $k$ un corps) -- mais alors il faut considérer $\mathrm{Dist}(\unsure{x|X},k)$ comme module topologique.

% loc 15–20, Re NB
\uline{Re NB} Surtout intéressant pour \unsure{monoïde}. En plus \deleted{\unsure{de|du}} \inserted{\unsure{du sous foncteur}} $G$, $\mathrm{dist}(G)$ \unsure{'} [\unsure{resp.} $\mathrm{Dist}(G,k)$] contient \unsure{toute} l'information infinitésimale de $G$ : p.ex. l'\uline{algèbre de Lie} de $G$ est \unsure{contenue|contenu} dans $\mathrm{Dist}(G\unsure{,|.}k)$, -- \deleted{\unsure{et les op.}} sa loi de $[\,]$ est induite par le crochet de l'algèbre, et en car. pos sa puiss. \unsure{pième} aussi -- de $=$ que l'hyperalgèbre.

\begin{center}\rule{1.6cm}{0.55pt}\end{center}
% loc 21, between the two rules
Cas des schémas constants : on retrouve l'\uline{algèbre} d'un groupe fini.

\begin{center}\rule{1.6cm}{0.55pt}\end{center}
% loc 22, second Cas
Cas du groupe diagonalisable
% loc 23, A = k(M)
\[
 A=k(M)\quad\text{algèbre du groupe à \unsure{coeff.|coeffs} dans } k
\]
% loc 24–25, two paired rows
\diagram{\begin{tikzcd}[ampersand replacement=\&,column sep=small,row sep=large]
 A\otimes A \arrow[r,dash] \& A \& \text{provient de} \& M\times M \arrow[r] \& M \& \text{loi de groupe}\\
 A\otimes A \& A \arrow[l] \& \text{---} \& M\times M \& M \arrow[l,swap,"\substack{\text{application}\\\text{diagonale}}"] \&
\end{tikzcd}}
% loc 26–27, lowest two lines
\[
\begin{aligned}
 &\mathrm{Dist}(\unsure{\mathbb D|D}(M)(k)=\deleted{\unsure{*}}k^{M}\quad\text{avec la loi d'\uline{algèbre} \uline{produit}}\\
 &\qquad\text{[Combinaisons \unsure{lin|lin.} infinies de \unsure{caractères}]} \qquad
 \text{(\unsure{ne dépend pas de la str. additive de|ne dépend pas de la str. d'addition de} $M$)}
\end{aligned}
\]
% archive sheet number: 22
